\documentclass[addpoints,spanish, 12pt,a4paper]{exam}
\pointpoints{punto}{puntos}
\hpword{Puntos:}
\vpword{Puntos:}
\htword{Total}
\vtword{Total}
\hsword{Resultado:}
\hqword{Ejercicio:}
\vqword{Ejercicio:}

%\printanswers

\usepackage[utf8]{inputenc}
\usepackage[spanish]{babel}
\usepackage{eurosym}
\usepackage{yhmath}
%\usepackage[spanish,es-lcroman, es-tabla, es-noshorthands]{babel}

\usepackage{verbatim}

\usepackage[margin=1in]{geometry}
\usepackage{amsmath,amssymb}
\usepackage{multicol}

\usepackage{graphicx}
\graphicspath{{../img/}} 

\newcommand{\class}{Matemáticas 4º Aplicadas}
\newcommand{\examdate}{\today}
\newcommand{\examnum}{Recuperación - Problemas Aritméticos}
\newcommand{\tipo}{D}


\newcommand{\timelimit}{50 minutos}

\renewcommand{\solutiontitle}{\noindent\textbf{Solución:}\enspace}

\pagestyle{head}
\firstpageheader{\includegraphics[width=0.2\columnwidth]{header_left}}{\textbf{Departamento de Matemáticas\linebreak \class}\linebreak \examnum}{\includegraphics[width=0.1\columnwidth]{header_right}}
\runningheader{\class}{\examnum}{Página \thepage\ of \numpages}
\runningheadrule

\pointsinrightmargin % Para poner las puntuaciones a la derecha. Se puede cambiar. Si se comenta, sale a la izquierda.
\extrawidth{-2.4cm} %Un poquito más de margen por si ponemos textos largos.
\marginpointname{ \emph{\points}}


\begin{document}

\noindent
%\begin{tabular*}{\textwidth}{l @{\extracolsep{\fill}} r @{\extracolsep{6pt}} }
%\textbf{Nombre:} \makebox[3.5in]{\hrulefill} & \textbf{Fecha:}\makebox[1in]{\hrulefill} \\
% & \\
%\textbf{Tiempo: \timelimit} & Tipo: \tipo 
%\end{tabular*}
%\rule[2ex]{\textwidth}{2pt}
%Esta prueba tiene \numquestions\ ejercicios. La puntuación máxima es de \numpoints. 
%La nota final de la prueba será la parte proporcional de la puntuación obtenida sobre la puntuación máxima. Para la evaluación de pendientes de 3ºESO o 2ºPMAR se tendrán en cuenta los apartados 1, 2, 3 y 6: 
%
%\begin{center}
%
%
%\addpoints
% %\gradetable[h][questions]
%	\pointtable[h][questions]
%\end{center}

\noindent
\rule[2ex]{\textwidth}{2pt}

\begin{questions}

\question Obtén el valor de x en cada caso:
\begin{parts}
\part  x = 2\% de 135

\begin{solution}$135\cdot0,02=2,7$\end{solution}
\part  28\% de x = 15,96

\begin{solution}15,96/0,28=57\end{solution}
\end{parts}
\addpoints


\question  El precio de un artículo sin IVA es de 315\euro. Si he pagado 365,40\euro, ¿qué porcentaje de IVA me han cargado?
\begin{solution}
Se ha pagado de IVA 365,40 - 315 = 50,40 \euro. 50,40/315 = 0,16. 16\%
\end{solution}

\question  He pagado 35,7\euro por una camisa que tenía un 15\% de rebaja. ¿Cuál era su precio antes de estar rebajada?
\begin{solution}
85\% de x = 35,7. 35,7/0,85=42\euro
\end{solution}

\question  Carlos coloca 18000\euro \   al 3,5\% anual y los mantiene en el banco durante 3 años, eligiendo la modalidad de interés simple. ¿A cuánto ascienden el capital obtenido durante los tres años?
\begin{solution}
18000x0,035=630. 630x3=1890. 19890\euro
\end{solution}


\question  Calcula en cuánto se transforman 9500\euro \ colocados al 3,5\% de interés compuesto anual durante 3 años.
\begin{solution}
$9500\cdot1,035^3=10532,82$\euro
\end{solution}

\question  Raquel, María e Isabel han ganado un premio de 8000\euro en un sorteo. Sabiendo que, para comprar los boletos, Raquel puso 5\euro, María 8\euro e Isabel 12\euro, ¿cuánto le corresponderá a cada una del premio que han ganado?
\begin{solution}
1600, 2560, 3840\euro
\end{solution}



\question Un juguete vale en una juguetería 40 euros. Durante las fiestas navideñas sube un 22\% y una
vez que éstas han pasado, baja un 9\%. Calcular su precio final y el porcentaje de subida
respecto al precio inicial
\begin{solution}
44.408 y 11.02\%
\end{solution}


\addpoints

\end{questions}

\end{document}
